\subsection{An operational interpretation of the tariff restriction}\label{sec:institution63}
One interpretation is a provider that has already accepted a service-credit obligation. Histories index observable eligibility classes; preliminary service and a randomized terminal credit jointly discharge the promised obligation. The expected cap limits the average obligation to a class, while the realization ceiling limits any delivered outcome. A command book is the set of terminal credit levels configured before individual assignments. A common opening charge represents a standardized setup procedure; exceptional levels require additional integration or qualification. This is a model interpretation, not a claim that a particular provider has supplied or adopted these inputs.

The accepted aggregate promise is a commitment to be met exactly. It is not obtained by changing an upper spending budget into an equality. The operational decision is therefore how to implement an already contracted obligation, including its fixed installation charges. The unrestricted model permits concave terminal benefits and costly preliminary service. The tariff theorem identifies a specific physical regime---a common linear benefit and zero incremental preliminary-service cost---in which fee structure alone changes computational tractability. Heterogeneous caps, ceilings, probabilities, catalog spacing, and an unrestricted command allowance remain present. Near-standard tariffs are treated by the explicit perturbation certificate in Section~\ref{sec:robust63}, rather than by silently rounding measured charges.
